% 本文件由 LaTeX 排版 Agent 根据有效 translation.json 自动生成。
% author=Tertullian; work_id=0404; title=致其妻

\chapter{致其妻}

% structure: p0001 source=h1 target=omitted reason=page-title-already-rendered-by-parent

% structure: p0002 source=h2 target=section reason=relative-heading-rank; base=h2

\section{卷一}

% structure: p0003 source=h3 target=subsection reason=relative-heading-rank; base=h2

\subsection{第一章 论文的宗旨。否认写作的个人动机}

我深思，我至爱的在主内的同仆，即使在这样早的时期，就应为你在我离世后——若我先你被召——所必须遵循的道路作出预备，并将这预备的遵守托付于你的诚信。因为在世俗的事上，我们足够活跃，且希望我们每个人的利益得到照顾。若我们为此\Emph{此类}事立遗嘱，何以不更应该为我们在神圣与天上的事上预先考虑我们的后代，并在某种意义上遗赠一份在遗产分配之前即可领受的遗赠——（我指的是）关于那些从\OriginalTerm{不朽的美物}{immortal goods}和天上的产业中所分配之物的劝诫与证明？唯一的是，愿你能够完整地接受我这劝诫的信托，但愿天主赐予；愿荣耀、光荣、声誉、尊威和权能归于他，从现在直到永远！\par

因此，我所给你的是这条诫命：尽你所能的一切恒心，在我们离世后，你不要再婚；并非你会因此赐予我什么益处，除非是你使自己获益。但对于基督徒，在他们离世之后，在复活的日子，没有复婚的应许，因为他们将被转化为天使的情状和圣洁。因此，由肉体的嫉妒所生的焦虑，在复活的日子，即使是在他们选择为代表的那位曾连续嫁给七位兄弟的妇人身上，也不会伤害她那么多丈夫中的任何一个；也没有任何丈夫等候她使她困惑。撒杜塞人所提出的问题已屈服于主的判决。不要以为我为了为自己保全你肉体至终的完全奉献，因我怀疑（预期）轻慢的痛苦，就在这样早的时期向你灌输（永久）寡居的劝告。在那日，我们之间将不再重新开始淫逸的耻辱。天主并未向他的（仆人）应许如此轻浮、如此不洁之事。但无论是对你，还是对任何其他属于天主的女人，我们所给的劝告都将有益，我们请允许详细讨论。\par

% structure: p0006 source=h3 target=subsection reason=relative-heading-rank; base=h2

\subsection{第二章 婚姻合法，但非多妻制}

我们并未禁止男人与女人的结合，这结合蒙天主祝福，作为人类的苗圃，为着充满大地和装备世界而被设计，因此是允许的，但只限单独一次。因为亚当是厄娃的惟一丈夫，厄娃是他惟一的妻子，一个女人，一根肋骨。我们承认，在我们的祖先和族长们当中，不仅结婚是合法的，甚至多妻也是合法的。那时也有妾侍。但虽然教会曾以比喻的方式出现在会堂中，然而（简单地解释）有必要设立（某些事物），这些事物日后应被剪除或修改。因为法律（届时）将随之而来。（这还不够：）因为补足法律不足的原因应当先行于（那位要补足这些不足者）。因此，紧接着法律，天主的圣言必须来临，引入属灵的割损。因此，借着那些日子的广泛许可，为后来的修正预先提供了材料，主藉他的福音，然后宗徒在（犹太）时代的末了，或剪除了多余的东西，或整理了混乱。\par

% structure: p0008 source=h3 target=subsection reason=relative-heading-rank; base=h2

\subsection{第三章 婚姻是好的：独身更可取}

但不要以为我如此详述古时所允许的自由与后来所加的限制，是为了教导说基督的来临是要解散婚姻，废除婚姻的束缚；仿佛从此以后我就要禁止结婚。让那些在他们的其他乖谬中教导分开\BibleQuote{二人成为一体}的人去负责吧；他们否认那位从男性取女性后，在婚姻的计算中，将出于同一物质团体中的两个身体结合于彼此之间的天主。总之，我们从未在任何地方读到婚姻被禁止；当然，因为它们是\Quote{一件好事}。然而，什么是比这\Quote{好}更\Emph{好}的，我们从宗徒那里学到，他确实\Emph{允许}结婚，但\Emph{更偏爱}禁欲；前者是因为诱惑的诡计，后者是因为时局的艰难。现在，藉着观察每个命题如此给出的理由，很容易看出准许结婚的权力是基于\Emph{需要}；但凡是\Emph{需要}所赐予的，凭其本性就贬低了它。事实上，经上写着：\BibleQuote{与其燃烧，不如结婚}，请问，这个\Quote{好}的性质是什么，它只是与\Quote{恶}相比较而被称赞，以致\Quote{结婚}之所以更\Emph{好}，仅仅是因为\Quote{燃烧}更\Emph{差}？不然！既结婚又不燃烧岂不是更好吗？为什么，即使在迫害中，利用所准许的权柄\BibleQuote{从这城逃到那城}也\Emph{更好}，胜过被捉住、受拷打而否认信德。因此，那些在蒙福的信仰宣认中能够离世的人更为有福。我可以说，所\Emph{准许}的不是\Emph{好}的。因为情况如何？如果我被捉住并宣认信仰，我必然要死。如果我认为那命运可叹，那么（逃跑）是好的；但如果我害怕那被准许的事情，那么这个准许的事情就有某种可疑之处。但那\Quote{更好}的，没有人曾\Quote{准许}它，因为它是不言自明的，并且其固有的纯洁性显而易见。有些东西之所以不可\Emph{渴望}，仅仅是因为它们未被\Emph{禁止}，尽管当别的东西比它们更受偏爱时，它们在某种意义上是被\Emph{禁止}的，因为对更高事物的偏爱是对最低事物的劝阻。一件事物并不只是因为不是\Quote{恶}就是\Quote{好}，也不只是因为不是\Quote{有害}就是\Quote{恶}。此外，那全然\Quote{好}的事物之所以优秀，是因为它不仅有百利而无一害，而且还有益。因为你必须选择有益的，而不是仅仅无害的。因为每场战斗的目标都是\Emph{第一}，\Emph{第二}则附有安慰，但不是胜利。但如果我们听从宗徒，忘记背后，努力向前，追求更好的赏报。这样，虽然他没有在我们身上\BibleQuote{撒下罗网}，但他指出了有益的事情，当他说：\BibleQuote{没有丈夫的妇人，挂念主的事，为叫身心圣洁；但那已出嫁的，挂念的是怎样悦乐丈夫。}但是，他从未准许结婚到不希望我们竭力效法他的榜样的地步。有福的人将如\Person{保禄}一样！\par

% structure: p0010 source=h3 target=subsection reason=relative-heading-rank; base=h2

\subsection{第四章 论血肉的软弱及类似托辞}

但我们在经上读到\BibleQuote{血肉是软弱的}；为此我们在某些情况下安慰自己。然而我们也读到\BibleQuote{心神是刚强的}；因为两个子句出现在同一句话中。血肉是属地的，心神是属天的材料。那么，为什么我们这些太容易自我原谅的人，只提出我们软弱的部分，却不看那刚强的部分呢？为什么属地的不应服从属天的呢？如果心神比血肉更强，因为它也有着更高贵的起源，那么我们跟随软弱的，就是自己的过错了。现在，有两种人性的软弱使得婚姻对那些脱离婚姻的人来说成为必要。第一种也是最有力的，是来自\Emph{肉欲}；第二种来自\Emph{世俗的贪欲}。但是，我们这些天主的仆人，既弃绝享乐又弃绝野心，应当拒绝这两者。肉欲声称成年期的功能，渴求美貌的丰收，以自己的羞耻为乐，为女性提出丈夫的必要性，作为权威和安慰的来源，或使其免于恶意的谣传。针对这些计谋，请你应用我们那些名字在主内的姐妹们的榜样：当她们的丈夫先她们而去时，她们不让任何美貌或年龄的机会优先于圣洁。她们宁愿与天主结婚。她们的美貌、她们的青春都奉献给天主。她们与祂同住，与祂交谈，日夜\BibleQuote{触摸}祂；她们将祈祷作为妆奁献给主；每当她们愿意，就从祂那里得到嘉许作为嫁妆。这样，她们为自己赢得了主的一项永恒恩赐；在地上，由于不结婚，她们已经被视为属于天使的家族。藉着这些女子的榜样，训练自己效法她们的恒心，你将藉着属灵的情感埋葬那肉欲，以不朽恩惠的补偿来消除美貌和青春的短暂易逝的欲望。\par

另一方面，我前面提到的这种\Emph{世俗的}\OriginalTerm{贪欲}{concupiscence}有其原因：荣耀、贪婪、野心、不足之感；这些原因编造出结婚的\Quote{需要}——自欺欺人地许诺将得到天上的事物作为回报——即，在别人的家庭中称霸，依仗别人的财富，从别人的积蓄中榨取奢华，挥霍你感受不到的支出！愿这一切远离信徒，他们并不挂虑生计，除非我们不信任天主的应许，不信任祂的眷顾和上智安排；祂以如此恩宠装饰田间的百合花；祂无需花儿的劳作就喂养天上的飞鸟；祂禁止我们为明天的食物和衣服忧虑，应许祂知道祂每个仆人所需的——其实不是沉重的项链，不是累赘的衣裳，不是高卢骡子或日耳曼挑夫，这些都增添婚礼的光彩；而是\Quote{\OriginalTerm{足用}{sufficiency}}，这适合节制和谦逊。我求你，你若\Quote{服事主}，就当认为自己一无所缺；不，你若拥有万有的主，你就拥有一切。常思念天上的事，你就会轻看地上的事。对于在主面前盖印密封的寡居生活，除了坚忍，别无所需。\par

% structure: p0013 source=h3 target=subsection reason=relative-heading-rank; base=h2

\subsection{第五章 论子嗣之爱作为婚姻的托词}

人们为自己提出的进一步结婚理由来自于对后裔的忧虑，以及生儿育女的苦涩、苦涩之乐。对我们而言，这是无谓的。因为我们为何渴望生育儿女呢？当我们有了他们，又渴望在他们之前（指进入荣耀）送走他们（我是说，鉴于当前即将来临的苦难）；我们自己也渴望被带离这个最邪恶的世界，被接入主面前，这甚至是宗徒的愿望！天主的仆人，难道真的需要后裔吗！我们对自己的救恩已经足够有把握，以致有闲暇为儿女操心！我们必须为自己寻求负担，而大多数外邦人尚且避免这些负担，他们（外邦人）受法律强制，因堕胎而人口减少；这些负担最终对我们最不相宜，因为它危及信德！因为主为何预言\Quote{怀胎的和乳养的有祸了}，除非祂见证，在那个卸下重担的日子，儿女的\OriginalTerm{负担}{encumbrances}会成为不便？当然，这些负担属于婚姻；但那（\Quote{祸}）不会属于寡妇。她们（\Emph{她们}）在天使的第一声号角中将无碍地跃出——将自由地忍受到底一切压力与迫害，腹中没有婚姻的沉重果实起伏，怀中也没有。\par

因此，不论是为了肉体、为了世界、还是为了后裔而结婚，这些\Quote{需要}无一能影响天主的仆人，以致阻止我认为一次永久地向其中某样屈服，并藉一次婚姻平息所有此类\OriginalTerm{贪欲}{concupiscence}就够了。让我们天天结婚，并在我们的结婚中让我们像索多玛和哈摩辣那样被那可怕的日子追上！因为那里不仅仅只有婚姻和买卖；但当祂说\Quote{他们又娶又买}时，祂给肉体和世界首要的\OriginalTerm{恶习}{vices}打上了烙印，这些恶习最使人远离神圣的训导——一是由于放纵宴乐，另一是由于贪图获取。然而，那\Quote{盲目}在\Emph{那时}早在\Quote{世界的终局}之前就被感受到了。那么，如今天主若阻止我们陷于\Emph{古时}在祂面前就已可憎的这些\OriginalTerm{恶习}{vices}，情况将如何呢？宗徒说：\Quote{时间是紧迫的。从此以后，那有妻子的要像没有妻子一样。}。\par

% structure: p0016 source=h3 target=subsection reason=relative-heading-rank; base=h2

\subsection{第六章 外邦人的榜样劝勉守寡与独身}

但是，如果那些\Emph{有}妻子的人尚且被要求忘却他们所有的，那么那些\Emph{没有}的人，岂不更应禁止再去寻求那已不再拥有的吗？如此，那丈夫已离世的女人，应从那时起藉着禁婚使她的性别得到安息——许多外邦妇女正是出于对爱侣的怀念而奉献此禁婚！当某事看似困难时，让我们看看那些应对更大困难的人。有多少人从受洗那一刻起，就在她们肉身上盖上了童贞的印记？又有多少人，通过平等的相互同意，解除了婚姻的债务——为了渴望天国而自愿成为阉者？然而，如果婚姻的纽带仍然完好时尚且能忍受禁欲，那么在它已经解开之后，岂不更应如此？因为我以为，完全舍弃那完好无损的，比不再渴望已失去的更为困难。诚然，一位圣洁的妇女在丈夫去世后为天主的缘故而贞洁，这已是足够艰难而劳苦的事，何况那些外邦人为了尊崇他们自己的撒殚，竟忍受涉及童贞和守寡的司祭职分！例如在罗马，那些与那\Quote{不灭之火}的预表打交道的人，与那古龙一起守护他们自己（未来）刑罚的征兆，他们便是基于\Emph{童贞}而被选立的。在\Place{埃吉翁}城，一位\Emph{童贞女}被分配给亚该亚的朱诺；而那些在\Place{德尔斐}发狂的（女祭司）不知婚姻。此外，我们知道\Emph{寡妇}侍奉非洲的色列斯；她们确实是被一种极其严厉的遗忘从婚姻中诱离：因为她们不仅离开尚在世的丈夫，甚至还为丈夫引进其他妻子来代替自己——丈夫们当然对此微笑——禁止她们与男性有任何接触，甚至包括亲吻自己的儿子；然而，她们凭着持久的操练，在这种守寡的纪律中坚持，甚至排除了圣洁情感的慰藉。这些规条是魔鬼赐给它的仆人的，而它竟被听从！它挑衅天主的仆人，仿佛以己方的贞洁与之同等！连地狱的司祭也是贞洁的！因为它找到了即使在善行中也毁灭人的方法；对它而言，用纵欲杀死一些人，用贞洁杀死另一些人，并无区别。\par

% structure: p0018 source=h3 target=subsection reason=relative-heading-rank; base=h2

\subsection{第七章 丈夫的死亡是天主对寡妇贞洁的召唤。来自圣经和异教的进一步证据}

对我们而言，贞洁已被救恩之主指出为获得永生的工具，并作为（我们）信德的见证；作为我们这肉身的赞许，这肉身将被滋养以迎接那即将来临的\Quote{不朽的衣裳}；最后，为承受天主的旨意。此外，我劝你思考：没有一个人是违背天主旨意而离世的，如果连一片叶子不经过天主也不会从树上落下。那带我们入世的，也必然要带我们离世。因此，当丈夫因天主的旨意去世时，婚姻也因天主的旨意而结束。为何\Emph{你}要恢复天主所终止的？为何你通过重复婚姻的奴役，拒绝那赐给你的自由？宗徒说：\BibleQuote{你已与妻子绑定，不要寻求解脱；你已从妻子解脱，不要寻求绑定}。因为即使你\BibleQuote{\Emph{犯罪}}再婚，他仍说\BibleQuote{肉身将有压力}。因此，尽力而为吧，让我们喜爱贞洁的机会；一旦它出现，让我们决心接受，以便我们在婚姻中未能坚持的，在守寡中得以坚持。必须抓住那终结了\Emph{必要性}所命令的机会。第二次婚姻对信德有何等损害，对圣洁何等阻碍，教会的纪律和宗徒的规条都已宣告，当时他不允许再婚者治理（教会），他不愿接纳寡妇进入那行列，除非她曾是\BibleQuote{一个男人的妻子}；因为天主的祭台应当保持纯洁。那环绕教会的整个光环代表（由）圣洁组成。司祭职是外邦人中的守寡和独身的（职务）。当然（这是）符合魔鬼竞争的原则。因为异教之王，大司祭，再婚是非法的。圣洁在天主眼中是何等悦乐，竟连祂的仇敌都效法它！——当然，不是因为它与任何善事有亲和力，而是作为倨傲地效法那悦乐主天主的事。\par

% structure: p0020 source=h3 target=subsection reason=relative-heading-rank; base=h2

\subsection{第八章 结论}

因为，论到守寡在天主眼中所享有的尊荣，借着先知，主的一句话中有一个简短的总结：\Quote{你们应公正对待寡妇和孤儿；然后来，让我们辩论，主说。} 这两个名称，托付给天主仁慈的照顾，既然她们缺少人助，万有的父便承担保护她们。请看，寡妇的恩人竟被置于与寡妇同等的地位，而她的捍卫者将\Quote{与主辩论！} 我认为，如此大的恩赐并未赐给贞女。虽然在\Emph{她们}身上，完美的贞洁和完整的圣洁将最接近天主的容颜，但寡妇的任务更为艰巨，因为不贪求你所不知道的，以及回避你从未后悔失去的，是容易的。更为光荣的是那意识到自己权利、知道自己所见之物的节德。贞女或许可以被认为更幸福，但\Emph{寡妇}的任务更艰难；前者因为她一直保持着\Quote{善}，后者因为她为自己找到了\Quote{善}。前者是恩宠加冕，后者是德行加冕。因为有些事物来自天主的慷慨，有些来自我们自己的努力。主所赐的恩惠由其恩宠调节；人努力追求的事物则通过热切追求而获得。因此，你要热切追求节德，它是端庄的使者；勤勉，它使女人不致成为\Quote{游荡者}；节俭，它轻视世俗。要追随配享天主的同伴和谈话，记住那被宗徒引用的短句：\Quote{不良的交往败坏良好的品德。} 多嘴、懒惰、贪杯、好奇的同伴，对寡妇的生活目标造成极大的伤害。通过多嘴，有损端庄的话语潜入；通过懒惰，引诱人偏离严格；通过贪杯，引入任何邪恶；通过好奇，灌输好色竞争的精神。这些女人中没有一人懂得谈论一夫之妻的好处；因为她们的\Quote{神}，如宗徒所说，\Quote{就是他们的肚腹}；同样，肚腹的邻居也是如此。\par

最亲爱的同仆，这些考虑，我如此早就向你推荐，虽然以宗徒之后来说是多余的处理，但对你可能是一种安慰，因为（如果情况如此）你将在其中珍视我的纪念。\par

% structure: p0023 source=h2 target=section reason=relative-heading-rank; base=h2

\section{第二卷}

% structure: p0024 source=h3 target=subsection reason=relative-heading-rank; base=h2

\subsection{第一章 写作这第二卷的原因}

最近，在主内最亲爱的同仆，我尽我所能，为你详细探讨了当婚姻（以任何方式）结束时，一位圣洁的女子应遵循何种道路。现在让我们关注次佳的忠告，关于人性的软弱；这受到某些例子的劝诫，当她们因离婚或丈夫去世而有机会实践节制时，不仅抛弃了获得如此大善的机会，甚至在再婚时也不愿遵守那条规则，即\Quote{要首先在主内结婚}。因此我的思绪陷入混乱，担心我自己劝你坚持一夫之妻和守寡，但现在可能因提及仓促的婚姻而给你造成\Quote{跌倒的机会}。但若你在智慧上是完美的，你当然知道哪条路更有用，你就必须持守哪条路。然而，由于那条路是艰难的，并非没有尴尬，因此它是（寡居）生活的最高目标，我在劝你时有所停顿；若不是因为我现在更加忧虑，也不会有任何理由再次向你提到这一点。因为肉体的节德越是高贵，服务于守寡，似乎越容易被原谅不去坚持。因为当事情困难时，宽恕就容易。但既然\Quote{在主内}结婚是允许的，因为是在我们的能力之内，那么\Emph{不}遵守你\Emph{能}遵守的，就更加有罪。此外，宗徒关于寡妇和未婚者，\Emph{劝告}她们永久保持那种状态，当他说：\Quote{但我愿意所有人都坚持（效法我的榜样）；} 但关于'在主内'结婚，他不再\Emph{劝告}，而是明确\Emph{命令}。因此在这件事上，如果我们不服从，我们就冒险，因为忽略一个\Quote{劝告}比忽略一个\Quote{命令}可能更不受惩罚；前者来自\Emph{建议}，向\Emph{意志}提出（接受或拒绝）；后者来自\Emph{权威}，并归于\Emph{必然}。在前者，忽视似乎是\Emph{自由}；在后者，是\Emph{悖逆}。\par

% structure: p0026 source=h3 target=subsection reason=relative-heading-rank; base=h2

\subsection{第二章 论宗徒在\BibleRef{格前 7:12-14}中的含义}

因此，当近日某妇人将其婚姻从教会界限中移出，而与外邦人结合时，我忆及此等事在以往已由他人所为：我对其自身的任性或顾问的两面手法感到惊奇，因无任何圣经许可此等行为——我说：\Quote{我惊奇，}我说，\Quote{他们是否以格林多前书的那段经文自慰，其中写道：\InnerQuote{若某弟兄有不信主的妻子，而妻子同意婚姻，就不应离弃她；同样，信主的妇人若嫁给不信主者，若丈夫同意（继续结合），也不应离弃他：因为不信主的丈夫因信主的妻子而圣洁，不信主的妻子因信主的丈夫而圣洁；否则你们的子女就不洁净。}}可能因普遍理解此劝言是针对已婚信者，他们以为（由此）甚至获准与不信者结婚。但愿如此解释（经文）者并非故意自陷网罗！但显然这段圣经指向那些因天主的恩宠而在外邦婚姻中被发现的信者；按经文本身所言：\Quote{若}它说，\Quote{某信者\Emph{有}不信主的妻子；}它并未说，\Quote{\Emph{娶}不信主的妻子。}它表明，已与不信主妇人生活在婚姻中、目前因天主的恩宠而皈依的人，有责任继续与妻子同在；诚然，为此缘故，以免任何人获得信德后，以为必须离弃如今在某种意义上为\Quote{外人}和\Quote{陌生人}的妇人。因此他随后还加了一个理由，即\Quote{我们\Emph{在平安中}蒙召归于主天主；}且\Quote{不信主者可能通过婚姻的使用，被信者\Emph{赢得}}。该段结尾完全确认（推断）应如此理解。\Quote{各人}它说，\Quote{如何蒙主召叫，就应如何持守。}但\Quote{蒙召}的是外邦人，我认为，而非信者。但若他（在此讨论的经文里）是绝对地就信者的婚姻发言，（那么）他（实际上）已给予圣徒随意结婚的许可。然而，若他给予如此许可，他绝不会随后加上一句如此与自己的许可相异相反的话，说：\Quote{妇人，丈夫死后，便自由了：她可随意嫁人，\Emph{只应在主内}。}在这里，无论如何，无需再考量；因为那里\Emph{可能}有考量之处，圣神已预言式地宣告了。为免我们误用他所说的话，\Quote{她可随意嫁人}，他加上了\Quote{只应在主内}，即，在主的名内，这无疑是指\Quote{与基督徒结婚}。因此，那宁愿寡妇和未婚女子持守贞洁的\Quote{圣神，}那劝勉我们效法祂自己的圣神，除\Quote{在主内}外，没有规定再婚的任何其他方式；唯独在此条件下，祂才容许放弃节欲。\Quote{只}他说，\Quote{在主内}；祂在祂的律法上加了一个重担——\Quote{\Emph{只}}。无论你用何种语气和方式说出这个词，它都是沉重的：它既命令又建议；既禁令又劝勉；既请求又威胁。它是简短精炼的一句话；且正因其简短而雄辩。天主的声音常如此（说话），使你立刻明白，立刻遵守。因为谁能不明白宗徒预见到此类婚姻对信德有许多危险和创伤，因而予以禁止？并且他首先防范圣洁的肉身在外邦肉身中被玷污。此时有人说：\Quote{那么，蒙主拣选归祂而在外邦婚姻中的人，与早（即婚前）已是信者的人有何区别，以致他们不可同样谨慎对待自己的肉身？——前者被禁止与不信者结婚，后者却被命令继续婚姻。为何，若我们被外邦人玷污，岂不该像不被捆绑者一样被分开？}我若蒙圣神赐予（能力），将回答；首先，（其他）所有论证之前，陈述主更喜悦婚姻不被缔结，胜于在一切情况下被解散：简短地说，祂禁止离婚，除非为奸淫的缘故；但祂称赞节欲。因此，前者有继续的必要；后者更有甚至不结婚的能力。其次，若按圣经，那些在外邦婚姻中被信德\Quote{捕获}的人，不因此被玷污，因为连同他们自己，他人也被圣化：无疑，那些在婚前已被圣化的人，若与\Quote{异教肉身}混杂，就不能圣化那（与他们）结合时未被\Quote{捕获}的（肉身）。此外，天主的恩宠圣化它所\Emph{找到}的。因此，未能被圣化的就是不洁的；不洁的与圣洁的没有分，除非以其本性玷污并摧毁它。\par

% structure: p0028 source=h3 target=subsection reason=relative-heading-rank; base=h2

\subsection{第三章。对前一章所提及的某些\Quote{危险与创伤}的评论}

若果真如此，那么信徒与外邦人缔结婚姻无疑是犯了奸淫，并应按照宗徒的书信，\BibleQuote{与这样的人连一同吃饭也不可}，而被排除于弟兄的共融之外。或者我们\Quote{到了那一天}，要在主的审判台前拿出婚约，声称祂所禁止的婚姻已正式缔结？此处所禁止的并非\Quote{通奸}，也非\Quote{奸淫}。接纳一个陌生男子（到你的卧榻）较少玷污\Quote{天主的圣殿}，较少将\Quote{基督的肢体}与娼妓的肢体相混合。据我所知，\BibleQuote{我们不属于自己，而是被重价买来的}；何等重价？天主的血。因此，伤害我们这肉身，就是直接伤害祂。那说\Quote{娶一个\InnerQuote{外人}固然是罪，却是小罪吗？}的人，究竟何意？而在其他情况下（撇开对属于主的肉身的伤害），\Emph{每一件}故意得罪主的罪都是大罪。因为，越是能够避免的，就越负有不服从的罪责。\par

现在，让我们数一数宗徒所预见到的信德的其他危险或创伤（如我所说）；这些危险不仅对肉身，同样对灵魂\Emph{也}极其严重。因为谁怀疑日常与不信者交往会逐渐抹去信德呢？\BibleQuote{恶语败坏善行}，何况是共同生活与不可分割的亲密关系！任何信主的女子都必须服从天主。然而，她怎能服侍两个主——主和她的丈夫，何况还是外邦人？因为服从外邦人，她就会行外邦人的事——追求个人魅力、装饰头发、世俗的优雅、低贱的谄媚，甚至婚姻的秘密也被玷污：不像在圣徒中，夫妻的本分以对必要性的尊重、以谦逊和节制、如同在天主眼前那样光荣地履行。\par

% structure: p0031 source=h3 target=subsection reason=relative-heading-rank; base=h2

\subsection{第四章 不信主的丈夫给妻子设置的障碍}

但她要注意如何向丈夫尽本分。无论如何，她不能按照纪律的要求使主满意；因为她身边有一个魔鬼的仆人，是他\Emph{自己的}主人的代理人，来阻碍信徒的追求和本分：因此，若要守斋戒站，丈夫天一亮就约妻子去浴室相见；若有斋戒要守，丈夫同一日会设宴；若要行慈善，家庭事务从未如此紧急。因为谁会容忍妻子为了探望弟兄，挨街串巷去别人家里，甚至去所有更贫穷的茅屋？谁愿意忍受她因夜间集会（如果需要的话）而离开他？最后，谁会毫无忧虑地容忍她在逾越节庆典整夜不归？谁会毫无怀疑地让她去参加那被他们诽谤的主的晚餐？谁会容许她溜进监狱亲吻殉道者的锁链？不，实际上，去会见任何一位弟兄行亲吻礼？为圣徒的脚预备洗脚水？从自己的食物和杯子中分给他们一些？渴慕他们？心中怀念他们？若有朝圣的弟兄到来，在异教徒的家中他如何得到款待？若要施舍给某人，谷仓和库房都被锁上了。\par

% structure: p0033 source=h3 target=subsection reason=relative-heading-rank; base=h2

\subsection{第五章 即使与\Quote{容忍的}丈夫一起也会遇到的罪与危险}

\Quote{但有些丈夫容忍我们的（做法），并不烦扰我们。}\Emph{这里}，因此便有了罪：因为外邦人知道了我们的（做法）；因为我们受制于不义之人的知情；因为多亏他们我们才能做任何善工。容忍一件事的人不能不知情；或者，若因他不容忍而瞒着他，他便是被惧怕的。但既然圣经命令两件事——即我们做工为主而不让任何第二人知道，且不使自己受压——那么你在哪一方面犯罪并无区别：无论是对付容忍你的丈夫的知情，还是因避免他的\Emph{不}容忍而自己受苦。祂说：\BibleQuote{不要把你们的珍珠投在猪前，免得它们用脚践踏，并且转过来咬伤你们。}\BibleQuote{你们的珍珠}便是你们日常言谈中的鲜明标记。你们越小心隐藏它们，就越会使之可疑，越易招致外邦人好奇的抓握。你岂能避人耳目地在床上、在身上画十字，吹走某些不洁，甚至夜间起来祈祷？难道不会被当作在行某种魔法？你的丈夫难道不知道你在进食前偷偷品尝的是什么？若他知道那是饼，他岂不以为就是那被说的饼？而每一个（丈夫），若不知这些事的缘由，能不抱怨、不怀疑那是饼还是毒药，就简单容忍吗？有些人，\Emph{的确}容忍；但那是为了践踏、嘲笑这样的女子；他们保留这些秘密，以备他们相信的危险，万一他们曾受伤害：他们容忍妻子，却用她们的嫁妆，以基督徒的名号来要挟，作为沉默的代价；同时威胁她们，要请某个间谍作仲裁的诉讼！大多数女子没有预见这一点，往往因此而发现自己的财产被勒索，或失去信德。\par

% structure: p0035 source=h3 target=subsection reason=relative-heading-rank; base=h2

\subsection{第六章 被迫参加异教仪式及宴乐的危险}

天主的女仆居于异族的劳役中；在这些劳役中，每逢所有魔鬼的纪念日，所有君王的庆典，年初，月初，她都会被馨香之气搅扰。她将不得不走出（家门），经过一个挂着月桂冠、悬着灯笼的大门，仿佛从某个新的公共情欲的议事厅出来；她将不得不常与丈夫一同坐在俱乐部中，常坐在酒馆中；她从前惯常服侍\Quote{圣人}，却有时不得不服侍\Quote{不义者}。她岂不会因此认出自己受罚的预判，因为她\Emph{服侍}那些她（从前）期待\Emph{审判}的人？她将渴望谁的手？她将分享谁的杯？她的丈夫将唱什么给她听，或她唱给丈夫听？我想，那以天主为晚餐的人，从酒馆中会听到什么！从地狱中，有什么关于天主的提及？有什么对基督的呼求？信德藉着圣经的穿插（在交谈中）的培育在哪里？圣神在哪里？慰藉在哪里？神圣的祝福在哪里？一切都是陌生的，一切是敌对的，一切是被定罪的；被恶者用来磨灭救恩！\par

% structure: p0037 source=h3 target=subsection reason=relative-heading-rank; base=h2

\subsection{第七章 与异教徒婚后妻子皈依的情况：截然不同，远更有希望}

如果这些事也可能发生在那些身处外邦人婚姻中、获得信德后仍继续此状态的妇女身上，她们仍可被谅解，因为是在这些环境中\Quote{被天主所擒获}；她们被\Emph{命令}坚守婚姻状态，且被圣化，并有希望\Quote{得利}。那么，如果这种（在皈依\Emph{之前}缔结的）婚姻在天主面前是有效的，为什么（皈依\Emph{之后}缔结的）婚姻就不能顺利进行，以致不受这些压力、困境、阻碍和玷污的折磨，既然它已（部分地）获得了神圣恩宠的认可？因为，一方面，前一种情况下的妻子，从外邦人中蒙召实践某种卓越的天上德行，藉着某种显著（神圣）眷顾的可见证据，令其外邦人丈夫畏惧，使他不太愿意烦扰她，不太积极设陷阱害她，不太勤于窥探她。他感受到了\Quote{大能作为}；他看到了经验性的证据；他知道她变好了：因此，连他自己也因恐惧而成为天主的候选人。这样，对于天主恩宠已建立亲密关系的人，更容易被\Quote{赢得}。但另一方面，不请自来地、自发地降入禁地，则是另一回事。不讨主喜悦的事，当然得罪主，当然是由恶者引入的。这事实的一个标志是：只有\Emph{求婚者}才觉得基督徒的名字令人愉悦；因此，有些异教徒男子竟然不厌恶躲避基督徒女子，正是为了消灭她们、夺走她们、将她们排除在信德之外。既然这种婚姻是由恶者谋取、却被天主定罪，你就有理由不必怀疑它绝不可能有好的结局。\par

% structure: p0039 source=h3 target=subsection reason=relative-heading-rank; base=h2

\subsection{第八章 从异教法律中得出的论据，反对与不信者通婚。结尾处详述信仰伴侣结合的幸福}

让我们进一步探究，仿佛我们真是神圣判决的审问者，看它们是否合法（如此被定罪）。即使在异教国家中，难道所有最严厉、最严守纪律的主人不是禁止自己的奴隶与外人通婚吗？——当然是为了使他们不致陷入放纵、荒废职责、将主人的财产供给外人。再者，难道（异教国家）没有规定，那些在主人正式警告后仍坚持与其他人的奴隶交往的女子，可以被当作奴隶索取？难道地上的纪律要比天上的诫命更严格，以至\Emph{外邦}女子若与外人结合就失去自由，而\Emph{我们的}女子与魔鬼的奴隶结合，却仍保持（原来的）地位？诚然，她们会否认主曾藉祂自己的宗徒给予她们任何正式警告！\par

除了信德的软弱，那常倾向于世俗欢乐的\OriginalTerm{贪欲}{concupiscentiae}，我还能将这种疯狂归咎于什么呢？这确实主要发生在较富裕的人中；因为一个人越富有，并因\Quote{贵妇}之名而自大，她就需要更大的房子来承载她的负担，就像一片田地，让野心可以驰骋。在这样的人看来，教堂显得寒酸。在天主的殿中，富人很难找到；若是在那里找到了，也很难找到未婚的。那么，她们该怎么办？除了从魔鬼那里，她们到哪里去找一个能供养她们轿辇、骡子和身形怪异的卷发器的丈夫呢？一个基督徒，即使富有，或许也供不起这些。我求你，将外邦人的榜样摆在自己面前。大多数外邦女子，出身高贵、家财万贯，毫无区别地与卑贱低微之人结合，为自己寻找（丈夫），或为享乐，或为纵欲而令其残缺。有些女子与自己的自由民和奴隶同居，不顾舆论，只要能找到不阻碍她们自由的丈夫。对于一位有信德的基督徒来说，与一位财产不如自己的信友结婚是令人不快的，因为她注定要在一个贫穷的丈夫身上增加自己的财富！因为如果\BibleQuote{穷人，而非富人，是天国的}，那么富人在穷人身上将找到更多（比她带给他的，或者比她在富人身上找到的）。她将从那位在神前富有的人那里得到更丰厚的嫁妆。让她在地上与他平等，而在天上或许不会如此。是否有必要怀疑、探究、反复思考，天主将祂自己的财物托付给他的人是否适合给予嫁妆？我们从何处能找到足够的言语来充分描述那婚姻的幸福？这婚姻由教会缔结，由祭献确认，由祝福标记并封印；天使将消息带回天上，圣父认为它有效。因为在地上，儿女没有父亲的同意，就不能合法地结婚。两个信友的结合是何等样式？他们同有一个希望、一个愿望、一个纪律、一个服务。两人同为弟兄，同为仆人，在神魂或肉身上没有区别；不，他们实在是\BibleQuote{两人成为一体}。哪里肉身是一个，神魂也是一个。他们一同祈祷，一同俯伏，一同守斋；互相教导，互相劝勉，互相扶持。两人同在天主的教会中，同在天主的筵席上，同在困境、迫害和安慰中。彼此不隐藏什么，彼此不回避，彼此不成为负担。病人得到探访，穷人得到周济，且毫无顾虑。施舍而无受折磨之虞，参与祭献而无疑虑，日常的虔诚不受阻碍：没有偷偷的画十字，没有颤抖的问候，没有无声的祝福。两人之间回响着圣咏和圣歌；他们互相挑战，看谁能更好地向主歌唱。基督看到、听到这些事，就欢喜。祂将自己的平安赐给这些人。哪里有两个（人），哪里也有祂自己。哪里有祂，哪里就没有恶者。\par

这些就是宗徒的那段话在其简洁之下留给我们去理解的事。如果有必要，将这些事提示给你自己的心。借着这些，使你自己远离某些人的榜样。\Emph{否则}结婚，对信友来说，不是\Quote{许可的}，也不是\Quote{有益的}。\par

\SourceRecord{Translated by S. Thelwall.；Ante-Nicene Fathers；Vol. 4.；Edited by Alexander Roberts, James Donaldson, and A. Cleveland Coxe.；Buffalo, NY: Christian Literature Publishing Co.,；1885.；<http://www.newadvent.org/fathers/0404.htm>.}
